\subsection{\texorpdfstring{\(\mathbf{R}^n\) 的加法群}{R 的 n 次幂的加法群}}

集合 \(R^{n}\) 赋以 \(R^{n}\) 的 n 个因子的加法群结构之积的群结构，是一个交换群；我们使用加法记号，\(x = (x_{i})\) 与 \(y = (y_{i})\) 的和因此是 \(x + y = (x_{i} + y_{i})\)。数空间的拓扑与这个群结构相容；赋以这两个结构，\(R^{n}\) 是一个拓扑群，称为 n 维实数空间的加法群。若 n = 0，我们约定 \(R^{0}\) 表示只由单位元素组成的群。

这个群的一致结构称为 \(\mathbf{R}^n\) 的加法一致结构，它是 \(\mathbf{R}^n\) 的诸因子的一致结构的积（第 III 章，§ 3，第 2 小节）。若对每个整数 \(p > 0\) ，\(\mathrm{V}_p\) 表示由点对 \((x, y)\) （\(\mathbf{R}^n\) 中的）组成的满足 \(\max |x_i - y_i| \leqslant 1/p\) 的集，则诸集 \(\mathrm{V}_p\) 构成一个基本

\[
\mathbf {1} \leqslant i \leqslant n
\]

围域系，即该一致结构的围域系。每当我们把 \(R^{n}\) 看作一致空间时，除非明确地声明相反的情形，我们心中所指的总是上面刚刚定义的加法一致结构。赋以这个一致结构，\(R^{n}\) 是完备一致空间（第 II 章，§ 3，第 5 小节，命题 10）。

